%=============================================================================!
% 5.E  EXERCISES                                                              !
%=============================================================================!
\unnumbered{Exercises}
\label{sec:ro-exercises}

The exercises run from questions of understanding, through calculations,
to short derivations marked `show that'. Full solutions are in
\hyperref[sec:ro-solutions]{`Solutions to the exercises'} at the end of
the book, and Notebook~05 checks every numerical answer.

\begin{exercise}[a door pushed at a slant]
\label{ex:ro-door}
Why are door handles placed far from the hinges? A push of
\SI{10}{\newton} at the handle, \SI{0.8}{\metre} from the hinge, is made at
\ang{30} to the face of the door instead of at right angles. What is its
torque about the hinge?
\end{exercise}

\begin{exercise}[pulling the puck further]
\label{ex:ro-pull}
The puck of \cref{sec:ro-plane} ($m = \SI{0.2}{\kilogram}$, circling at
\SI{0.5}{\metre} and \SI{1.2}{\metre\per\second}) is pulled in to
\SI{0.1}{\metre}. Find its speed of going round, its kinetic energy and the
work done by the hand, neglecting the slow inward motion.
\end{exercise}

\begin{exercise}[a straight line]
\label{ex:ro-straight}
Show that a body moving at constant velocity, with no force on it, has a constant
angular momentum $m\,(\vb{r} - \vb{a})\times\vb{v}$ about any point
$\vb{a}$, its position measured from $\vb{a}$, and explain, with the triangles
swept in equal times, why it sweeps equal areas in equal times.
\end{exercise}

\begin{exercise}[a cross product]
\label{ex:ro-product}
Find $(2, 0, 1)\times(1, 3, -1)$, check that it is at right angles to both
vectors, and find the area of the parallelogram they span in two ways.
\end{exercise}

\begin{exercise}[motion in a plane]
\label{ex:ro-plane}
A body moves with a constant angular momentum $\vb{L} \neq \vb{0}$ about the
origin. Show that it stays in the plane through the origin at right angles
to $\vb{L}$.
\end{exercise}

\begin{exercise}[changing the point]
\label{ex:ro-origin}
Show that the angular momentum of a group of bodies about the point
$\vb{a}$, $\sum_i(\vb{r}_i - \vb{a})\times\vb{p}_i$, equals $\vb{L} -
\vb{a}\times\vb{P}$, with $\vb{L}$ its angular momentum about the origin
and $\vb{P}$ its total momentum. When is it the same about every point?
\end{exercise}

\begin{exercise}[a rod]
\label{ex:ro-rod}
Show that a thin uniform rod of mass $M$ and length $\ell$ has the moment
of inertia $M\ell^2/12$ about an axis through its middle at right angles to
it, and $M\ell^2/3$ about a parallel axis through one end.
\end{exercise}

\begin{exercise}[the parallel axis]
\label{ex:ro-parallel}
A body has the moment of inertia $I_{\mathrm{cm}}$ about an axis through
its centre of mass. Show that about a parallel axis a distance $s$ away its
moment of inertia is $I_{\mathrm{cm}} + Ms^2$, and check the result
against \cref{ex:ro-rod}.
\end{exercise}

\begin{exercise}[the stool again]
\label{ex:ro-stool}
The stool of \cref{sec:ro-rigid}, now with $I_0 =
\SI{2}{\kilogram\metre\squared}$ and weights of \SI{3}{\kilogram} held at
\SI{0.7}{\metre}, turns at \SI{1.5}{\radian\per\second}. The weights are
pulled in to \SI{0.3}{\metre}. Find the new angular speed and the work done
by the arms.
\end{exercise}

\begin{exercise}[a rolling ball]
\label{ex:ro-sphere}
A solid ball has $I = \frac25 MR^2$ about any axis through its centre.
Rolling from rest, how fast is it moving after a drop of \SI{1}{\metre}, and
how long does it take to roll \SI{2}{\metre} down a slope of \ang{30}?
Where does it finish in the race of \cref{fig:ro-rolling}?
\end{exercise}

\begin{exercise}[the grip a disc needs]
\label{ex:ro-friction}
Using Newton's second law for the centre of mass along the slope, find the
friction force that the slope must supply to a disc of \SI{1}{\kilogram}
rolling down at \ang{30}, and show that in general it is $\kappa Mg\sin
\alpha/(1 + \kappa)$.
\end{exercise}

\begin{exercise}[a square of atoms]
\label{ex:ro-square}
Four atoms of \SI{1}{\amu} sit at $(\pm1, \pm1, 0)$\,\si{\angstrom}. Find
their inertia tensor about the origin and its principal moments. Why is
every axis through the origin in the plane of the square a principal axis?
\end{exercise}

\begin{exercise}[a flat molecule]
\label{ex:ro-flat}
Show that for a body whose parts all lie in the plane $z = 0$, $I_{zz} =
I_{xx} + I_{yy}$, and that $I_{xz} = I_{yz} = 0$, so that the $z$ axis is a
principal axis.
\end{exercise}

\begin{exercise}[a diatomic molecule]
\label{ex:ro-diatomic}
Two atoms of masses $m_1$ and $m_2$ lie on the $z$ axis a distance $d$
apart. Show that their inertia tensor about their centre of mass is
diagonal, with entries $m_{\mathrm r}d^2$, $m_{\mathrm r}d^2$ and 0.
\end{exercise}

\begin{exercise}[elimination]
\label{ex:ro-solve}
Solve, by elimination, $4\omega_x + \omega_y = 6$, $\omega_x + 4\omega_y +
\omega_z = 12$, $\omega_y + 4\omega_z = 14$.
\end{exercise}

\begin{exercise}[what is left]
\label{ex:ro-left}
Two atoms of \SI{1}{\amu} at $(\mp1, 0, 0)$\,\si{\angstrom} move with
velocities $(0.01, 0.01, 0)$ and $(-0.01, 0.03,
0)$\,\si{\angstrom\per\femto\second}. Remove the drift and the turning,
describe the motion that remains, check \cref{eq:ro-split} in
\si{\amu\angstrom\squared\per\femto\second\squared}, and give the total
kinetic energy in \si{\electronvolt}.
\end{exercise}

\begin{exercise}[counting]
\label{ex:ro-count}
How many internal degrees of freedom have ammonia, NH$_3$, and hydrogen
peroxide, H$_2$O$_2$, neither of them linear?
\end{exercise}

\begin{exercise}[a zero mode by hand]
\label{ex:ro-null}
A single spring of stiffness $k$ joins an atom at the origin to one at $(0,
0, d)$. Write down its $6\times6$ Hessian, and show directly that it turns
the displacements of a small turning about the $x$ axis into zero. Why does
a turning about the $z$ axis not count?
\end{exercise}

\begin{exercise}[a spinning bond stretches]
\label{ex:ro-stretch}
The separation of a turning LiF molecule goes round a circle, which needs a
pull $m_{\mathrm r}\omega^2 r$ towards the centre, supplied by the bond as
a spring, $k(r - d)$. Show that the bond stretches by about $d\,(\omega/
\omega_0)^2$, with $\omega_0 = \sqrt{k/m_{\mathrm r}}$
the angular frequency of the vibration, and evaluate it for the turning of
\cref{sec:ro-molecules}.
\end{exercise}

\begin{exercise}[a hotter molecule]
\label{ex:ro-hot}
At \SI{1000}{\kelvin} a LiF molecule carries on average $\kB T =
\SI{0.08617}{\electronvolt}$ in its turning. How long does one turn take?
\end{exercise}
